\documentclass{article}
\usepackage[UTF8]{ctex}
\usepackage{amsmath,amssymb,amsfonts}
\usepackage{enumitem}
\usepackage{geometry}
\geometry{a4paper, margin=1in}
% === 错题本排版优化（不参与题目标识） ===
\setlength{\parindent}{0pt}
\setlength{\parskip}{0.35em plus 0.10em minus 0.08em}
\setlength{\emergencystretch}{3em}
\setlist{itemsep=0.25em, topsep=0.40em, parsep=0pt}
\allowdisplaybreaks[2]
% === 错题本排版优化结束 ===
\begin{document}
\section*{北京邮电大学2026年数学分析真题}
\section*{15．北京邮电大学2026年数学分析试卷}

\subsection{题目1}
一、求极限 \(\lim _{n \rightarrow \infty} \sum_{i=1}^{n} 2^{\frac{i}{n}} \frac{1}{n+\frac{1}{i}}\) 。

\subsection{题目2}
二．求极限 \(\lim _{x \rightarrow 0} \frac{\int_{0}^{x} e^{\frac{t^{2}}{2}} \cos t \mathrm{~d} t-x}{\left(e^{x}-1\right)^{2}(1-\cos x) \arctan x}\) ．

\subsection{题目3}
三．求函数 \(f(x, y)=\left(y+\frac{x}{3}\right) e^{x+y}\) 的极值．

\subsection{题目4}
四．设 \(f(x)\) 在 \(x=0\) 处存在二阶导数，且 \(\lim _{x \rightarrow 0} \frac{f(x)}{x}=0\) ，证明：\(\sum_{n=1}^{\infty}\left|f\left(\frac{1}{n}\right)\right|\) 收敛。

\subsection{题目5}
五．已知平面区域 \(D=\left\{(x, y) \mid x^{2}+y^{2} \leq 2 y\right\}\) ，计算二重积分 \(\iint_{D}(x+y)^{2} \mathrm{~d} x \mathrm{~d} y\) ．

\subsection{题目6}
六．计算第二型曲线积分

\[
I=\int_{L} x y \mathrm{~d} x+(x-y) \mathrm{d} y+x^{2} \mathrm{~d} z .
\]

其中 \(L\) 为螺旋线：\(x=a \cos t, y=a \sin t, z=b t\) ，从 \(t=0\) 到 \(t=\pi, a, b\) 为常数．

\subsection{题目7}
七．求 \(I=\iint_{S}\left(x^{2}+y^{2}+2\right) \mathrm{d} S\) ，其中 \(S\) 为球面 \(x^{2}+y^{2}+z^{2}=2 a x, a>0\) ．

\subsection{题目8}
八．设 \(a_{n}>0, n=1,2,3, \cdots\) ，若 \(\lim _{n \rightarrow \infty} \frac{a_{n+1}}{a_{n}}=l, l\) 为常数，证明： \(\lim _{n \rightarrow \infty} \sqrt[n]{a_{n}}=l\) ．

\subsection{题目9}
九．若函数 \(f(x)\) 在 \([1,+\infty)\) 上满足利普希茨条件，即存在某一个常数 \(L>0\) ，使对于任意的 \(x_{1}, x_{2} \in [1,+\infty)\) ，有 \(\left|f\left(x_{1}\right)-f\left(x_{2}\right)\right| \leq L\left|x_{1}-x_{2}\right|\) ，证明：\(\frac{f(x)}{x}\) 在 \([1,+\infty)\) 上一致连续．

\subsection{题目10}
十．设 \(f(x)\) 在 \((0,+\infty)\) 上满足 \(f(x)=f\left(x^{2}\right)\) ，且 \(\lim _{x \rightarrow 0^{+}} f(x)=\lim _{x \rightarrow+\infty} f(x)=f(1)\) ，证明：

\[
f(x) \equiv f(1), x \in(0,+\infty)
\]

\subsection{题目11}
十一。设 \(f(x)\) 在 \((-\delta, \delta)\) 上有 \(n+1\) 阶导数，且 \(f^{\prime}(0)=f^{\prime \prime}(0)=\cdots=f^{(n)}(0)=0, f^{(n+1)}(0) \neq 0\) ，由微分中值定理，对 \(|x|<\delta\) ，存在 \(\theta \in(0,1)\) ，使得 \(f(x)-f(0)=f^{\prime}(\theta x) x\) ，证明： \(\lim _{x \rightarrow 0} \theta=\frac{1}{\sqrt[n]{n+1}}\) ．

\subsection{题目12}
十二．设 \(f(x)\) 在 \([a, b]\) 上有二阶连续导数且 \(f^{\prime \prime}(x)>0\) ，证明：

\[
f\left(\frac{a+b}{2}\right) \leq \frac{1}{b-a} \int_{a}^{b} f(x) \mathrm{d} x \leq \frac{f(a)+f(b)}{2}
\]

\subsection{题目13}
十三。设二元函数 \(f(x, y)=|x-y| \varphi(x, y)\) ，其中 \(\varphi(x, y)\) 在点 \((0,0)\) 的一个邻域内连续，证明：函数\(f(x, y)\) 在 \((0,0)\) 处可微的充要条件是 \(\varphi(0,0)=0\) ．

\subsection{题目14}
十四．设函数 \(f_{n}(x)=x n^{k} e^{-n x}\) ，其中 \(k\) 为常数，讨论 \(\left\{f_{n}(x)\right\}\) 在 \([0,+\infty)\) 上的一致收敛性．

\subsection{题目15}
十五．设函数 \(f(x)\) 在区间 \([a, b]\) 上无界，证明：

\begin{itemize}
  \item[（1）] 存在 \(\left\{x_{n}\right\} \subset[0,1]\) ，使得 \(\lim _{n \rightarrow \infty} f\left(x_{n}\right)=\infty\) ．
  \item[（2）] 存在 \(c \in(a, b)\) ，使得对任意的 \(\delta>0, f(x)\) 在 \((c-\delta, c+\delta) \cap[a, b]\) 无界。
\end{itemize}
\end{document}
